% =====================================================================
\chapter{Comprehension: Inverse Compression}\label{ch:comprehension}
% =====================================================================

This chapter addresses the first half of Problem \ref{prob:mind}: for a large system that
cannot be read through from the micro level, at what level should it be understood, and how
does comprehension lead into development. The central claim is that comprehension is inverse
compression, whose optimal depth is given by the stationarity condition of a rate--distortion
objective and whose feasibility is decided by observability rank.

\section{Failure of Sequential Reading}\label{sec:why-sequential}

A natural way to read is to follow the physical order and the control flow of the source code:

\begin{quote}
\code{main()} $\to$ \code{module\_a} $\to$ \code{module\_b} $\to$
\code{function\_c} $\to$ \code{helper\_d} $\to$ \code{allocator} $\to$
\code{operating\_system} $\to \cdots$
\end{quote}

One then keeps following downwards and ends in infinite recursion. The failure is not a lack
of patience but structural:

\begin{principle}[Calling relations are not cognitive dependencies]\label{pr:call-vs-cog}
The function-call relation is the topology of \emph{implementation}; the cognitive-dependency
relation is the topology of \emph{comprehension}. The two are not isomorphic. In studying
transaction isolation one may need to know how the transaction manager uses locks and
snapshots, but not necessarily understand the entire virtual-memory subsystem of the operating
system; conversely, in investigating anomalous memory usage the high-level business semantics
of the transaction manager may be temporarily irrelevant.
\end{principle}

The more fundamental reason concerns the relation between the direction of generation and the
direction of comprehension. A system is designed \emph{top-down} (intent first, then
architecture, and only then the unfolding of code); the inverse process should therefore take
the same direction. Reading code bottom-up is like understanding a cat by counting all its
molecules: feasible in principle, absurd in practice.

The concession ``feasible in principle'' has a practical consequence: sequential reading is
not \emph{impossible} but \emph{dominated}. Proposition \ref{prop:stopping} gives its precise
form: beyond the stationary point, description length increases while task loss no longer
falls.

\section{Cognitive Maps}\label{sec:maps}

Confronted with a large monolithic repository, one should build a conceptual map before a file
map. The division of labour is:

\begin{itemize}
\item a file map answers ``where the code is'';
\item a conceptual map answers ``what the system is made of, why these parts exist, who
      depends on whom, and what rules allow them to collaborate''.
\end{itemize}

\begin{definition}[The four cognitive maps]\label{def:fourmaps}
For a system $S$ and a task $T$, the four cognitive maps are:
\begin{enumerate}[label=(\roman*)]
\item \term{system panorama}: compresses $S$ into $5$--$15$ major responsibility domains and
      answers: which responsibility domain is the most important? Why does each
      responsibility domain exist? What are the most important dependency relations among
      them? Which mechanisms are shared by several responsibility domains? Which states and
      behaviours decide the correctness of the whole system?
\item \term{key subsystem map}: for each key subsystem, answers eight questions:
      responsibility; boundary (what it is responsible for and what it is not); state (what
      state it owns and who may modify it); mechanism; interface; constraints; failure
      behaviour; and implementation entry points (which source files and tests can verify the
      above conclusions).
\item \term{key execution path map}: converts the static architecture into a model of dynamic
      behaviour. For example, the real questions behind a query path ``input $\to$ parsing
      $\to$ planning $\to$ execution $\to$ storage access $\to$ return'': when is the query
      plan created? Where is state kept during execution? Who manages the transaction
      snapshot? How is an error propagated upwards when a storage access fails? Which steps
      hold locks? Which operations may trigger disk I/O? Which resources must be released
      when a query is cancelled?
\item \term{contract map}: records the relations among components by each of the five edge
      classes of Definition \ref{def:relgraph}, in particular the \emph{oblivious set} of each
      layer, that is, what that layer is allowed never to care about.
\end{enumerate}
\end{definition}

Map (iv) is this paper's addition to conventional architecture documentation. Its necessity
follows directly from Proposition \ref{prop:oblivious}: a contract that has not declared its
oblivious set has not declared its compression boundary, and hence one cannot judge how far
the level above may rely on it.

\section{Operational Criteria for Comprehension}\label{sec:five-capabilities}

``Comprehension'' must be operationalised, or it cannot be decided whether comprehension has
been achieved. This paper adopts a five-capability criterion.

\begin{definition}[Criteria for comprehension]\label{def:understand}
An agent (human or machine) is said to \term{understand} a subsystem $M$ to the degree
required by a task $T$ if and only if it can do the following:
\begin{enumerate}[label=(U\arabic*)]
\item \term{Explain}: state why this subsystem exists;
\item \term{Predict}: predict what behaviour a given input, state change or error will cause;
\item \term{Locate}: decide which components to inspect when a fault occurs;
\item \term{Modify}: find the correct extension points and the constraints that must be
      respected when a feature is to be added;
\item \term{Verify}: know which tests or implementation details can substantiate its own
      judgement.
\end{enumerate}
This paper calls a state that attains (U1)--(U5) \term{operational abstract comprehension}.
\end{definition}

The relation between Definition \ref{def:understand} and Definition \ref{def:sufficient} is
direct: (U2) is the behavioural statement of $T$-sufficiency; (U1), (U3) and (U4) are concrete
query families in $T$; (U5) is the requirement that one can trace back to lower-level evidence,
namely the two clauses ``sufficiently accurate'' and ``verifiable'' of Definition
\ref{def:mdl}.

One diagnostic corollary matters here: if one can only recite the architecture diagram but
cannot predict what a change will affect, the current abstraction is not yet sufficient. This
and the seam detectors of \S\ref{sec:seams} are one thing seen from two directions.

\section{Reading Order}\label{sec:protocol}

\begin{principle}[Descent protocol]\label{pr:protocol}
Form the model top-down, descend to verify on demand, and never read sequentially.
\end{principle}

The steps are as follows.

\paragraph{The Discipline Layer.}
Read the README, architecture documents, architecture decision records, design papers and the
kernel documentation tree. About $1\%$ of the time establishes about $80\%$ of the skeleton
model. The output is a hand-drawn organ diagram: which large blocks the system has, what each
is responsible for, and how data flows.

\paragraph{Runtime Observation.}
Run the system first, read it afterwards. Runtime behaviour exposes structure and is far more
honest than static code: observe the process and thread model; set breakpoints or trace the
whole journey of one typical request; read logs and metrics, \emph{the vocabulary of the logs
is the system's self-description}. One end-to-end trace is worth ten thousand lines of static
reading.

\paragraph{Interface Boundaries.}
Read the header files, public interfaces and exported symbols of modules in dependency order,
not their implementations. An interface is a compression boundary (Proposition
\ref{prop:oblivious}) and states what this layer promises; how it discharges that promise is
not yet needed. What is wanted at this point is an organ atlas, not a tissue section.

\paragraph{Data Shape.}
The data model is the most stable abstraction and control flow the most volatile. Establish
first: what the core entities are, who owns their lifecycle, which state machines exist, and
what is writable and what is read-only. Once the shape of the data is in hand, control flow is
only data rearranged over time.

\paragraph{Descent on Demand.}
Descend into code only with a concrete task in hand: fix this defect, add this feature.
Reading is then directed retrieval rather than wandering. Descent follows the V-shape ``call
site $\to$ implementation $\to$ underlying mechanism'': stop as soon as the current problem
can be explained, then return to the upper layer and reassess the model.

\paragraph{Prediction and Verification.}
At each layer, write the mental model down first (``I guess the executor fetches the next
tuple like this''), then verify it. \emph{Being able to predict behaviour is comprehension}.
If the prediction is right, the compression model is effective; if it is wrong, that spot is
exactly where the system is most counter-intuitive.

\paragraph{Write-Back to the Contract Map.}
Write the newly acquired knowledge back into the contract map of Definition \ref{def:fourmaps},
filling in that layer's new guarantees, new assumptions and new oblivious set. This step is
the only means of preventing cognitive drift: without write-back, the layered model gradually
drifts away from the real system.

\begin{remark}[An operational ratio]
Understand about $10\%$ of the code in depth while knowing where the other $90\%$ is and who
owns it. The ratio is not a precise target but a quantified reminder of the principle that
comprehension is not exhaustion.
\end{remark}

\section{Detection Criteria for Module Boundaries}\label{sec:seams}

Architecture is ``the places where a change is localised''. The key operation of comprehension
is therefore \emph{finding seams}. The criteria below are ordered by reliability, and all can
be executed mechanically or semi-mechanically.

\begin{enumerate}[label=(S\arabic*)]
\item \term{Where types stop crossing}. Where is data re-encoded, wrapped or degraded to a
      generic type (byte arrays, strings, maps)? This is a strong signal of a module boundary:
      where the type system stops providing guarantees is where a human must supply the
      contract.
\item \term{Units that can be tested in isolation}. Where does a test suite exist that can be
      run without starting the whole system? This is the sharpest criterion: \emph{only a unit
      that can be verified in isolation is a real module}. Its reliability comes from
      Definition \ref{def:sufficient}: a unit test is essentially the assertion that ``the
      interface observation is a sufficient statistic of behaviour''.
\item \term{Serialisation, version fields and copy boundaries}. Where are there version
      fields, serialisation, deep copies, cross-process or cross-network calls? Crossing such
      a boundary makes a contract visible and falsifiable.
\item \term{Centralised enforcement sites of invariants}. Where are there assertions, database
      constraints, type-level constraints or contract checks? Where invariants are enforced
      centrally is usually where a level boundary lies.
\item \term{Independent targets of the build system}. Where does the build configuration
      define an independent compilation target or package? This is a \emph{necessary} but not a
      \emph{sufficient} condition for a boundary: an independent target is not necessarily a
      real module, but a real module nearly always has one.
\item \term{Co-change graphs and incident history}. This is the most underrated technique.
      Examine the commit history: which files tend to be modified together in the same change?
      \emph{A co-change graph gives semantic coupling, whereas an import graph gives syntactic
      coupling}. Where the two disagree is exactly where the architecture documentation
      disagrees with reality. Incident and defect histories then give the real invariant gaps.
\end{enumerate}

\begin{principle}[Summary of the seam criteria]\label{pr:seam-summary}
A module is real if and only if its interface observation is a sufficient statistic of its
external behaviour; a boundary is real if and only if, once it has been crossed, changes on
the other side no longer need to be understood.
\end{principle}

Chapter \ref{ch:design} inverts Principle \ref{pr:seam-summary} and uses it as a design
objective.

\section{Active Identification and Passive Reading}\label{sec:active}

\begin{principle}[Comprehension as system identification]\label{pr:sysid}
Comprehending a complex system is essentially \emph{system identification}: constructing a
model that predicts its behaviour and quantifying the uncertainty of that model. The most
efficient form of system identification is the ``perturb--observe'' loop.
\end{principle}

Passive reading is an extremely poor way to sample, because it \emph{has no error signal}.
Every judgement it produces is unfalsifiable, and therefore cannot be updated. Active probing
is the opposite: form a falsifiable prediction first (``if I change this, only these two
modules should be affected''), make a minimal change or run a test, then observe the
prediction error. \emph{The error is the part of the structure you cannot see.}

\begin{proposition}[The advantage of active probing]\label{prop:active}
Let comprehension be an iterative update of a model $M$, in which each observation $o$
supplies information $I(\text{structure}; o \mid M)$ about the true structure. The observation
distribution of passive reading does not depend on $M$, so its expected information gain is
constant in $M$; the observation distribution of active probing depends on the current
prediction $\hat{y}(M)$, and its expected information gain is proportional to the prediction
error $\mathbb{E}\,|y - \hat{y}(M)|$. Hence, wherever the prediction error is non-zero, active
probing has strictly higher marginal information gain.
\end{proposition}

\begin{proof}[Argument]
The samples drawn by passive reading are independent of $M$, hence $p(o \mid M) = p(o)$ and
the information gain is constant. Active probing selects the observations that maximise
predictive disagreement (D-optimal or uncertainty sampling in experimental design), so
$p(o \mid M)$ concentrates on the region where the current predictions of $M$ are uncertain;
by the definition of information gain $I = \mathbb{E}[\log p(o\mid M) - \log p(o)]$, that
quantity is non-negative and vanishes when $p(o\mid M) = p(o)$, that is, when prediction and
truth agree. Hence active probing is strictly better wherever the prediction error is non-zero.
\end{proof}

For AI the conclusion carries more weight: the advantage of AI is not that it can ``read'' more
code but that it can \emph{probe} at high throughput, generating hypotheses, running tests,
making minimal changes, inspecting co-changes, reading incident records, and then updating the
model. That is a frequency humans cannot match.

\section{The Observability Criterion}\label{sec:observability}

\S\ref{sec:comprehension-formal} defines comprehension as task-constrained compression, and
Proposition \ref{prop:stopping} gives the stopping point of descent. Both presuppose that
continued observation can improve the model. This section identifies a stronger limitation:
\emph{when observation is insufficient, increasing the number of observations changes
nothing}.

\begin{definition}[Observation operator and task subspace]\label{def:observability}
Let $\Beh$ be the behaviour space of the system and $\mathcal{T} \subseteq \Beh$ the
\emph{task-relevant difference subspace}, that is, the subspace spanned by the differences
that the query family $Q$ of Definition \ref{def:task} can distinguish. Let
$\mathcal{J} : \Beh \to \mathcal{Y}$ be the observation operator \emph{available} to the
agent: it maps behaviour to what the agent can actually read (source code, logs, traces, test
results, runtime metrics). Define the \term{observability rank} on $\mathcal{T}$ as
\begin{equation}\label{eq:obsrank}
\mathrm{rank}_{\mathcal{T}}\mathcal{J} \;=\;
\dim \mathcal{J}\bigl(\mathcal{T}\bigr),
\end{equation}
the dimension of the task-relevant differences that remain distinguishable after
$\mathcal{J}$. The construction is isomorphic to the observability criterion of linear system
theory \cite{kalman1960}.
\end{definition}

\begin{proposition}[The observability criterion]\label{prop:observability}
An agent can identify the system to $T$-sufficiency (Definition \ref{def:sufficient}) if and
only if
\begin{equation}\label{eq:obscond}
\mathrm{rank}_{\mathcal{T}}\mathcal{J} \;=\; \dim \mathcal{T}.
\end{equation}
When \eqref{eq:obscond} fails, the kernel of $\mathcal{J}$ on $\mathcal{T}$ is non-trivial:
there exist two executions $e_1, e_2$ whose available observations are all identical while
their task-relevant behaviour differs. In that case \emph{no number of additional observations
can reduce task loss}, because new observations still lie in the range of the same
$\mathcal{J}$.
\end{proposition}

\begin{proof}
If \eqref{eq:obscond} holds, then $\mathcal{J}$ is injective on $\mathcal{T}$, task-relevant
differences are uniquely recoverable from observations, and $T$-sufficiency follows from
Definition \ref{def:sufficient}. If it fails, then
$\Ker \mathcal{J} \cap \mathcal{T} \ne \{0\}$; take a pair of executions $e_1 \ne e_2$
corresponding to a non-zero element of that kernel (in the task-relevant sense). The two are
indistinguishable under $\mathcal{J}$, so every $\mathcal{J}$-measurable decision rule gives
the same answer and is wrong on at least one of them. Adding observations does not change
$\mathcal{J}$, and therefore does not change that kernel.
\end{proof}

Proposition \ref{prop:observability} separates ``insufficient comprehension'' into two cases
that are entirely different in engineering terms:

\begin{principle}[Distinguishing two kinds of comprehension failure]\label{pr:two-failures}
\begin{enumerate}[label=(\alph*)]
\item \term{rank deficiency}: the observation operator itself cannot distinguish task-relevant
      differences. The only remedy is to \emph{change the observation}: add invariant checks,
      add trace points, add assertions, expose internal state. This is a design action, not a
      reading action.
\item \term{full rank but insufficient samples}: observation suffices to distinguish, but has
      not yet been collected. The remedy is to \emph{probe more}, with its efficiency governed
      by Proposition \ref{prop:active}.
\end{enumerate}
Confusing the two is the main reason advice such as ``read more code'' fails: when the problem
belongs to (a), additional reading is information-theoretically ineffective.
\end{principle}

\begin{remark}[Measuring the cost of comprehension correctly]\label{rmk:sample-complexity}
The cost of case (b) is a \emph{sample complexity}, not an amount of reading. If the contract
hypotheses are confined to a hypothesis class $\mathcal{H}$ (the set of operations in
Definition \ref{def:operad} provides a natural candidate), then the number of probes required
to identify the system to $T$-sufficiency at error $\varepsilon$ has a lower bound decided by
the distinguishing power of $\mathcal{H}$ on $T$, typically of the form
$\Omega(d/\varepsilon)$, where $d$ is the capacity dimension of the hypothesis class.

This accounting yields a testable prediction: \emph{the factor by which designed probing beats
exhaustive reading is quantifiable}, and that factor grows with the capacity of the hypothesis
class. Hypothesis $H_5$ of Chapter \ref{ch:program} thereby acquires a quantitative form.
\end{remark}

\begin{example}[A typical shape of rank deficiency]\label{ex:rank-deficient}
A service keeps all its state in process memory and exposes only the return values of requests.
If the task is to decide ``whether two concurrent requests can read each other's partial
updates'', then no number of return-value observations suffices, because the difference lies
entirely outside the return values. The correct action is not to read more code but to
introduce an observable: a runtime assertion of a concurrency invariant, or a test with
interleaving injection. The example shows that the distinction of Principle
\ref{pr:two-failures} is practical rather than merely taxonomic.
\end{example}

\section{Budget and Stopping Conditions}\label{sec:budget}

\S\ref{sec:comprehension-formal} has given the formal stopping criterion (Proposition
\ref{prop:stopping}). This section gives its engineering version and the allocation of budget.

\begin{principle}[Allocation of the descent budget]\label{pr:budget}
Allocate the comprehension budget in order of information gain per unit cost, not by code
volume or directory order. The typical high-yield entry points, in order, are: external
interface definitions and the data model $>$ one trace of an end-to-end execution path $>$
invariant enforcement sites $>$ co-change hotspots $>$ specific implementation details.
Low-yield entry points are: sequential reading of implementation bodies, reading through a
utility function library, and reading through generated code.
\end{principle}

\begin{principle}[Stopping conditions]\label{pr:stop}
Stop descending as soon as any of the following holds:
\begin{enumerate}[label=(\alph*)]
\item you can name every invariant governing the current change;
\item further descent no longer changes your prediction of ``what this change will affect'';
\item the cost of obtaining the next piece of evidence exceeds the expected rework cost of the
      current change.
\end{enumerate}
\end{principle}

Condition (a) corresponds to Principle \ref{pr:descent}; condition (b) to Proposition
\ref{prop:stopping}; condition (c) is an economic judgement: comprehension is not a matter of
more being better, but of \emph{enough}.

\section{Implications for Machine Comprehension}\label{sec:ai-comprehension}

The difficulty facing AI is not only a limited context length. Even when a model can read a
large amount of code, three distinct problems remain: too much information drowns the
important information; local pieces of information depend on one another in complex ways; and
comprehension of details does not automatically form a correct system-level model. Enlarging
the context window alone therefore does not solve every problem.

\begin{principle}[A software cognitive architecture for AI]\label{pr:ai-cog}
An effective machine cognitive architecture should have four layers:
\begin{enumerate}[label=(\roman*)]
\item \term{global cognitive model}: system goals, architecture diagram, key invariants, main
      data flows, core design decisions;
\item \term{subsystem cognitive model}: component responsibilities, interfaces, state
      ownership, dependency relations, lifecycle and error semantics;
\item \term{implementation-level index}: key functions, type definitions, call chains, source
      locations, tests and runtime evidence;
\item \term{verification and feedback}: compilation, tests, static analysis, execution traces,
      profiling, used to correct the models above.
\end{enumerate}
Its key property is: \emph{an upper layer model need not replicate all lower-level information,
but need only retain enough information for the agent to find the correct lower-level evidence
when it is required.}
\end{principle}

\begin{definition}[Traceable subsystem cognitive record]\label{def:record}
The cognitive record of each subsystem should contain at least nine fields, listed in Table
\ref{tab:record}.
\end{definition}

\begin{table}[htbp]
\centering
\small
\caption{The minimum field set of a subsystem cognitive record}\label{tab:record}
\begin{tabular}{L{3.0cm}L{11.4cm}}
\toprule
\textbf{Field} & \textbf{Role} \\
\midrule
Purpose & Why this subsystem exists \\
Contract & What guarantees it provides upwards (that is, $\Obs_i(\Beh_{i+1})$) \\
Oblivious & What this layer is allowed never to care about (that is, $\Obl(\Lambda_i)$) \\
State & What state it owns and who may modify it \\
Dependencies & Which other components it depends on \\
Invariants & Which conditions must never be violated, and where they are enforced \\
Lifecycle & How key objects are created, run and destroyed \\
Execution paths & How important requests or tasks pass through the system \\
Evidence & Which source files, tests or experiments support the above conclusions \\
Uncertainty & What has not been verified, and which assumptions may be wrong \\
\bottomrule
\end{tabular}
\end{table}

The last two entries of Table \ref{tab:record}, together with the Oblivious entry, are
additions of this paper, for the following reasons.

\emph{On Uncertainty}: AI cannot treat an architecture summary it has generated itself as
fact. It needs to know which conclusions come from actual source code, which come from tests,
and which are only provisional inferences. This is the same requirement as the epistemic
discipline of Chapter \ref{ch:objections}.

\emph{On Oblivious}: as stated above, it is the compression boundary itself. Without this
entry, a cognitive record describes only ``what there is'' and not ``what may be
disregarded'', and so cannot support the modification capability (U4).

\emph{On dynamic update}: when the code changes, the cognitive record must be updated in step.
A usable machine comprehension system should therefore not be a static knowledge base but a
\emph{bidirectionally traceable} relation between the source code and the abstract model.

\begin{principle}[Readability turns from a virtue into infrastructure]\label{pr:readability}
The ``comprehension interface'' of a system is becoming a first-class engineering deliverable.
Documentation used to be an appendix written for human readers; now machine-readable
architecture descriptions, the practice of interface-as-documentation, and agent-facing
working conventions are becoming components that a system must design, because the readership
of code now includes a non-human intelligence that relies only on the compression layer.
\end{principle}

The empirical content of this chapter is Hypothesis $H_5$ (Chapter \ref{ch:program}): a
comprehension process that uses a layered model together with active probing achieves a higher
task success rate per unit cost than exhaustive reading. Propositions \ref{prop:stopping},
\ref{prop:active} and \ref{prop:observability} hold under the stated model; Definition
\ref{def:understand} is a working definition whose value depends on whether it can be used to
distinguish ``comprehension'' from ``recitation'', which requires experimental test.
